\section{Background and Related Work}\label{sec:related}
\label{sec:related-serving}
\subsection{Serving mechanisms and workload variation}
Prefill processes an input context; autoregressive decoding generates new
tokens while retaining KV state. Iteration-level scheduling in Orca
\citep{yu2022orca}, PagedAttention in vLLM \citep{kwon2023vllm} and
stall-free chunked prefill in Sarathi-Serve \citep{agrawal2024sarathi}
address different batching and memory constraints. DistServe and Splitwise
separate prefill and decode resources \citep{zhong2024distserve,patel2024splitwise};
Llumnix rebalances instances \citep{sun2024llumnix}; Mooncake organizes
serving around KV state \citep{qin2025mooncake}. Fairness-aware token
accounting \citep{sheng2024vtc}, service-length-aware packing
\citep{jin2023s3}, learned ordering \citep{efficientllmscheduling2024}
and SLO-aware scheduling \citep{slai2025} illustrate why a single scalar
ranking need not preserve every deployment objective. We compare historical
simulator policies inspired by such mechanisms, not these systems as
interchangeable production implementations.

\label{sec:related-workloads}
BurstGPT, Azure and Bailian expose different production-derived request
populations \citep{burstgpt2025,azure2024trace,bailianqwen2026}.
Characterization and generation work examines workload heterogeneity,
including KV reuse and changes over time
\citep{kvcache_wild2025,ayearinllmserving2026,servegen2026}.
BurstGPT and Azure are separately released sources associated with Azure
infrastructure; their comparison does not isolate a provider effect.
Observed prompt lengths can also have different conversation/context
semantics. LSSP retains these distinctions when interpreting portability.

\subsection{Simulation and evaluation methodology}\label{sec:related-simulators}
Vidur and LLMServingSim provide infrastructure for evaluating serving
configurations \citep{vidur2024,llmservingsim2024,llmservingsim2_2025}.
LSSP does not claim a more accurate simulator. Its emphasis is the
comparative conclusion induced by a fixed model, together with the
restrictions on transporting that conclusion beyond the model.

\label{sec:related-dependence}
Workload-dependent performance is established background. Work on serving
workload selection \citep{papaioannou2024workload}, production-trace
dispatching \citep{yildiz2026dispatching}, and workload resampling in parallel
job scheduling \citep{zakay2014resampling} provides close methodological
precedents. The distinctive scope pursued here is a common, frozen panel
with jointly reported source, region and metric comparisons, explicit
reversal denominators, and uncertainty conditional on the available windows.
It is not a priority claim that workload sensitivity or ranking comparison
was previously unknown.

\label{sec:related-methodology}\label{sec:related-position}
Kendall's tau-b measures ordinal concordance while accommodating ties;
Spearman correlation and top-$k$ overlap provide complementary summaries.
Neither a positive aggregate correlation nor a shared top choice rules out
pair-specific reversals. Practical margins distinguish a numerically changed
order from an operationally meaningful difference under a stated convention.
These standard measures serve a reproducible comparison protocol. Our
finite-corpus recovery study estimates stability relative to the frozen
reference ranking, not a theorem about population sample complexity.

\subsection{What is added beyond standard ranking statistics}
Workload resampling has a direct performance-evaluation precedent.
Zakay and Feitelson's conference study decomposes traces into user
subtraces and recombines them while preserving activity cycles
\citep{zakay2013resampling}. The shared concern is that the structure of
the workload affects the reliability of a scheduler comparison. LSSP
resamples already selected request windows instead; it does not implement
user-level trace reconstruction or establish that its windows are
independent. The distinction prevents borrowing a validity argument from
a resampling design with different observational units.

Dem\v{s}ar's comparison framework uses ranks across multiple datasets to
compare algorithms, including Friedman omnibus tests and subsequent
comparisons \citep{demsar2006comparisons}. LSSP applies standard rank tools
to fixed policy implementations, but asks an additional descriptive
question: whether their order is preserved when the condition or metric
changes. A pooled difference between policies is not itself evidence of
such a change. A method can outperform another in every source and still
produce a strong omnibus statistic. The source-pair reversal table is
therefore a distinct object, not an automatic interpretation of Friedman
significance.

The methodological addition is the explicit linkage between comparison
objects: each ranking carries its implementation panel, observed-window
support, utility definition, practical threshold and tie rule. The
post-hoc audit shows that these choices have measurable consequences in
one corpus. For example, agreement of deterministic top choices and
intersection of maximizing sets yield different answers without any change
to the outcomes. This is evidence for reporting both objects when ties
are common; it is not a claim that a new estimator has been invented.

The study also separates three validation tasks that can otherwise be
conflated: reproducing an analysis from frozen numbers, assessing its
sampling assumptions, and checking whether the simulator represents the
intended serving mechanisms. Exact numerical reproduction addresses only
the first. The inferential audit addresses limits of the second; the
legacy native-runner audit exposes unresolved parts of the third. This
separation is the paper's relevance to modeling and performance evaluation,
but it also limits the strength of its systems conclusions. A benchmark
with correctly implemented standard statistics still needs a defensible
comparison model and sufficiently representative inputs.
